\documentclass[11pt,a4paper]{article}
\usepackage{amsmath,amssymb,amsthm}
\usepackage[margin=2.5cm]{geometry}
\usepackage{hyperref}

\newtheorem{definition}{Definition}[section]
\newtheorem{theorem}[definition]{Theorem}
\newtheorem{proposition}[definition]{Proposition}
\newtheorem{lemma}[definition]{Lemma}
\newtheorem{corollary}[definition]{Corollary}
\newtheorem{conjecture}[definition]{Conjecture}
\newtheorem{remark}[definition]{Remark}

\title{Hyperseed Formalization:\\
  Where Energy and AI Wars Collide}
\author{ProtomegaTron (automated formalization)}
\date{2026-09-18}

\begin{document}
\maketitle

\begin{abstract}
We formalize the intellectual structure of Ben Goertzel's Substack article
``Where Energy and AI Wars Collide'' (2026-03-03) within the Hyperseed
ontology. The article analyzes the convergence of energy infrastructure
and AI development as a geopolitical nexus. Key mappings: energy as base
resource for the cognitive fiber, energy bottleneck as base constraint on
fiber, centralization as energy-concentrated base, decentralization as
distributed base, energy topology determining AI topology, efficiency as
fiber-base coupling tightness, and the self-reinforcing cycle as fiber-base
positive feedback.
\end{abstract}

\tableofcontents

%% =================================================================
\section{Base resource: energy as fiber prerequisite}
\label{sec:base}
%% =================================================================

\begin{theorem}[Energy as base resource --- claims 1--4, 24--25]
\label{thm:base}
AI training and inference require enormous and growing amounts of energy.
Energy availability is becoming the primary bottleneck for AI scaling ---
more limiting than algorithmic innovation, data availability, or hardware
design. AI energy consumption is growing exponentially, driven by both
larger models and wider deployment.

In Hyperseed terms, energy is the \emph{base resource} that the cognitive
fiber requires. No energy, no computation, no fiber. This is the most
concrete fiber-base constraint:
\[
  |F| \leq \phi(E_{\text{base}}) \qquad \text{(fiber capacity bounded
  by base energy)}
\]

where $|F|$ is the fiber's computational capacity and $\phi$ is the
energy-to-computation conversion function (determined by hardware
efficiency).

This is a \emph{hard physical constraint}, unlike the softer constraints
of algorithmic design or data availability. You can design a more
efficient algorithm (improving $\phi$) or find more data (expanding fiber
content), but you cannot compute without energy (the base must exist).

The distinction between training and inference maps to a fiber-base
temporal distinction:
\begin{itemize}
\item \textbf{Training} = fiber construction: one-time, energy-intensive
  process that builds the fiber architecture. High base-energy demand,
  concentrated in time.
\item \textbf{Inference} = fiber operation: continuous, less intensive
  per query but growing in total. Ongoing base-energy demand, distributed
  over time.
\end{itemize}

As AI deployment scales, inference energy may dominate training energy ---
the fiber's operational energy cost exceeds its construction energy cost.
This parallels biological systems: growing a brain is metabolically
expensive, but running a brain over a lifetime costs far more total energy.
\end{theorem}

%% =================================================================
\section{Base topology determines fiber topology}
\label{sec:topology}
%% =================================================================

\begin{theorem}[Energy topology --- claims 5--9, 14--17, 26--28]
\label{thm:topology}
The topology of the energy base determines the topology of the AI fiber
that can exist over it. The fiber can't have a topology that the base
can't support.

\textbf{Centralization = energy-concentrated base:}
When the base (energy infrastructure) is concentrated in a few
locations/entities, the fiber (AI development) must be concentrated
too --- the fiber follows the base:
\[
  \text{Concentrated } E_{\text{base}} \implies
  \text{Concentrated } F \implies \text{Concentrated power}
\]

Only entities with massive energy access (big tech companies,
resource-rich nations) can train frontier models. The energy requirement
creates a massive barrier to entry, concentrating AI power in the hands
of those who control energy.

\textbf{Decentralization = distributed base:}
Distributed energy sources (solar, wind, local grids, micro-nuclear)
could support distributed AI development --- many small fibers rather
than one large one:
\[
  \text{Distributed } E_{\text{base}} \implies
  \text{Distributed } F \implies \text{Distributed power}
\]

If you want decentralized AI, you need decentralized energy first. The
energy topology is the prerequisite for the AI topology.

\textbf{Geopolitical implications:}
Control of energy infrastructure becomes a geopolitical lever for
controlling AI development. Compute energy is becoming a strategic
resource like oil was in the 20th century. Energy sovereignty (the
ability to power AI infrastructure independently) becomes essential
for AI sovereignty:
\[
  \text{Energy sovereignty} \implies \text{AI sovereignty}
  \implies \text{technological sovereignty}
\]

Energy access can be used as sanctions: restricting energy access to
a nation effectively restricts their AI capabilities. Energy wars and
AI wars converge because the same base resource (energy) constrains
both domains.

In Hyperseed terms, the energy base has geopolitical fiber: control of
the base determines control of the fiber, and geopolitical competition
for the base (energy) is simultaneously competition for the fiber (AI).
The base and fiber are not independent --- they form a coupled system
where base topology constrains fiber topology.
\end{theorem}

%% =================================================================
\section{Fiber-base coupling and feedback}
\label{sec:coupling}
%% =================================================================

\begin{proposition}[Efficiency and feedback --- claims 10--13, 22--23, 29--31]
\label{prop:coupling}
\textbf{Fiber-base coupling tightness:}
More efficient AI architectures have tighter fiber-base coupling: they
extract more cognitive work per unit of base energy:
\[
  \eta = \frac{\text{cognitive output}}{\text{energy input}}
  \qquad \text{(fiber-base coupling efficiency)}
\]

Hardware efficiency improvements (better chips, more efficient
architectures) increase $\eta$, partially offsetting growing demands.
Hyperon's hybrid architecture may be more energy-efficient than pure
transformer scaling --- combining neural, symbolic, and evolutionary
components that each operate efficiently in their respective domains,
rather than applying one energy-intensive approach to everything.

Alternative approaches (neuromorphic computing, hybrid architectures)
become more valuable as energy constraints bite --- they offer better
fiber-base coupling when the base is constrained.

\textbf{Self-reinforcing cycle = fiber-base positive feedback:}
The cycle of energy $\to$ AI $\to$ revenue $\to$ more energy is a
fiber-base positive feedback loop:
\[
  E_{\text{base}} \xrightarrow{\text{enables}} F
  \xrightarrow{\text{generates}} \text{value}
  \xrightarrow{\text{acquires}} E_{\text{base}}^+
  \xrightarrow{\text{enables}} F^+
\]

The fiber generates value that expands the base, which supports a
larger fiber, which generates more value. This is unstable --- it
drives toward monopoly unless counteracted by distributed base topology.

In competitive dynamics, the entity with the initial energy advantage
builds better AI, generates more revenue, acquires more energy, and
the advantage compounds. This is a classic positive feedback loop that
produces winner-take-all dynamics.

The countermeasure is distributed base topology: if many entities have
independent energy access (distributed solar, local grids, SMRs), the
positive feedback loop operates for many entities simultaneously rather
than concentrating in one. Distributed base prevents monopolistic fiber.
\end{proposition}

%% =================================================================
\section{Nuclear as high-density base point}
\label{sec:nuclear}
%% =================================================================

\begin{definition}[Nuclear renaissance --- claims 18--20, 29]
\label{def:nuclear}
Nuclear energy is experiencing renewed interest specifically because of
AI energy demands. Tech companies (Microsoft, Google, Amazon) are
investing in nuclear to secure compute energy.

In Hyperseed terms, nuclear provides a \emph{high-density base point}:
a large amount of energy concentrated in one location. This supports
large, concentrated fibers (big training runs that require sustained
high power).

Small modular reactors (SMRs) are particularly interesting because they
distribute these high-density points: instead of one massive power
plant, many smaller reactors can be deployed near compute infrastructure.
This creates a distributed-but-dense base topology:
\[
  \text{SMR array} = \text{distributed high-density base}
  \implies \text{distributed large-fiber topology}
\]

AI infrastructure needs reliable baseload power, not intermittent sources.
Nuclear provides this; solar/wind alone don't (without massive storage).
The base must be \emph{reliable} as well as sufficient --- intermittent
energy means intermittent computation, which is incompatible with
large-scale training runs that must run continuously for weeks.

The nuclear renaissance illustrates the base-determines-fiber principle:
the AI industry's need for a specific base topology (reliable, high-density,
ideally distributed) is driving investment in energy infrastructure that
provides exactly that topology.
\end{definition}

%% =================================================================
\section{Cross-references}
%% =================================================================

\begin{remark}[Connections to other formalizations]
\begin{itemize}
\item \textbf{Three Things the World Doesn't Understand} (2026-03-26):
  centralized vs.\ decentralized AGI. The energy topology argument
  provides the physical foundation: decentralized AGI requires
  decentralized energy.
\item \textbf{The Geopolitics of the Great AI Bet} (2026-06-19):
  geopolitics. The energy-AI convergence is a key component of the
  broader geopolitics of AI development.
\item \textbf{Big Tech's Big Bets on the Singularity} (2026-07-23):
  corporate strategy. Tech companies' nuclear investments are part of
  their bet on AI scaling --- securing the base resource for future
  fiber expansion.
\item \textbf{How Decentralized AI Can Help Avert} (2026-06-19):
  decentralization. Energy decentralization as an enabler of AI
  decentralization.
\item \textbf{Zuck's Version of the Future} (2026-07-14): corporate AI.
  Meta's AI strategy depends on energy access --- the base constraint
  shapes the corporate fiber.
\item \textbf{Beyond Human Comparison} (2026-03-19): efficiency.
  Energy efficiency maps to the Resourcefulness dimension of the
  Four-Factor Model --- how efficiently the system uses base resources.
\end{itemize}
\end{remark}

\begin{conjecture}[Energy topology as AGI predictor]
Conjecture: the topology of global energy infrastructure is a better
predictor of which entities will achieve AGI first than the topology of
their algorithms, data, or talent.

If this conjecture holds, the most important AGI race variable is not
``who has the best algorithm'' or ``who has the most data'' but ``who
has the most reliable, abundant, and efficiently coupled energy
infrastructure.'' The algorithm race is secondary to the energy race.

This suggests a policy implication: if you want to ensure that AGI is
developed by diverse, decentralized actors rather than monopolized by
a few energy-rich entities, the most effective intervention is not
AI regulation but energy infrastructure policy --- ensuring distributed
access to abundant, reliable energy. Energy policy is AI policy.
\end{conjecture}

\end{document}
